\documentclass{amsbook}

\usepackage{enumitem}

\newtheorem{theorem}{Theorem}[chapter]
\newtheorem{lemma}[theorem]{Lemma}

\theoremstyle{definition}
\newtheorem{definition}[theorem]{Definition}
\newtheorem{example}[theorem]{Example}
\newtheorem{xca}[theorem]{Exercise}

\theoremstyle{remark}
\newtheorem{remark}[theorem]{Remark}

\numberwithin{section}{chapter}
\numberwithin{equation}{chapter}

\begin{document}

\frontmatter

\title{Minqi's Solutions to the Exercises of Russell, Norvig (2009), Artificial Intelligence: A Modern Approach (3rd Edition)}

\author{Minqi Pan}

\date{\today}

\maketitle

\setcounter{page}{4}

\tableofcontents

\mainmatter

\chapter{Introduction}

\section{Exercise 1.1}
Define in your own words:
\begin{enumerate}[label=(\alph*)]
\item intelligence,
\item artificial intelligence,
\item agent,
\item rationality,
\item logical reasoning.
\end{enumerate}
\begin{proof}[Answer]
\begin{enumerate}[label=(\alph*)]
\item Intelligence is the ability to perceive information, perform reasoning and thinkings, retain knowledge, and take according actions.
\item Artificial intelligence is the study of the design of intelligent agents.
\item An agent is an independent entity that operates autonomously, perceives its environment, persists knowledge over a prolonged time period, adapts to changes, and creates and purses goals.
\item Rationality is a characteristic of an agent's behavior, that is, to always achieve the best outcome or the best expected outcome.
\item Logical reasoning is to reason according to the formal rules and norms of logic.
\end{enumerate}

\end{proof}
\end{document}
